% ============================================================
% CRYPTOGRAPHY TOOLKIT
% Theorem environments, algorithm layout, notation macros and
% TikZ styles for two-party protocol diagrams.
% This file is loaded at the end of packages.tex.
% ============================================================

% ---------- Packages ----------
\usepackage{amsthm}
\usepackage{mathtools}
\usepackage{algorithm}
\usepackage{algpseudocode}

\usetikzlibrary{arrows.meta,positioning,calc}

% ============================================================
% THEOREM ENVIRONMENTS
% All share one counter, numbered within the section
% (Definition 3.1, Theorem 3.2, ...), which is the usual
% convention in cryptography texts.
% ============================================================

\theoremstyle{definition}
\newtheorem{definition}{Definition}[section]
\newtheorem{construction}[definition]{Construction}
\newtheorem{example}[definition]{Example}
\newtheorem{remark}[definition]{Remark}

\theoremstyle{plain}
\newtheorem{theorem}[definition]{Theorem}
\newtheorem{lemma}[definition]{Lemma}
\newtheorem{corollary}[definition]{Corollary}
\newtheorem{proposition}[definition]{Proposition}
\newtheorem{claim}[definition]{Claim}

% ---------- cleveref names ----------
\crefname{definition}{Definition}{Definitions}
\Crefname{definition}{Definition}{Definitions}
\crefname{theorem}{Theorem}{Theorems}
\Crefname{theorem}{Theorem}{Theorems}
\crefname{lemma}{Lemma}{Lemmas}
\Crefname{lemma}{Lemma}{Lemmas}
\crefname{corollary}{Corollary}{Corollaries}
\Crefname{corollary}{Corollary}{Corollaries}
\crefname{construction}{Construction}{Constructions}
\Crefname{construction}{Construction}{Constructions}
\crefname{algorithm}{Algorithm}{Algorithms}
\Crefname{algorithm}{Algorithm}{Algorithms}
\crefname{lstlisting}{Listing}{Listings}
\Crefname{lstlisting}{Listing}{Listings}

% ============================================================
% ALGORITHM LAYOUT
% ============================================================

\algrenewcommand\algorithmicrequire{\textbf{Input:}}
\algrenewcommand\algorithmicensure{\textbf{Output:}}
\algrenewcommand\algorithmicforall{\textbf{for each}}

% Vertical guide lines for nested blocks
\algrenewcommand\algorithmicindent{1.1em}
\makeatletter
\renewcommand{\ALG@name}{Algorithm}
\makeatother

% ============================================================
% NOTATION — NUMBER THEORY
% ============================================================

\DeclareMathOperator{\lcm}{lcm}
\DeclareMathOperator{\ord}{ord}
\DeclareMathOperator{\Adv}{Adv}
\DeclareMathOperator{\poly}{poly}
\DeclareMathOperator{\negl}{negl}
\DeclareMathOperator{\Supp}{Supp}

% Sets and structures
\newcommand{\Zset}{\mathbb{Z}}
\newcommand{\Nset}{\mathbb{N}}
\newcommand{\Rset}{\mathbb{R}}
\newcommand{\Fset}{\mathbb{F}}
\newcommand{\Zn}[1]{\mathbb{Z}_{#1}}          % \Zn{n}  -> Z_n
\newcommand{\Zns}[1]{\mathbb{Z}_{#1}^{*}}     % \Zns{p} -> Z_p^*
\newcommand{\Fq}[1]{\mathbb{F}_{#1}}          % \Fq{q}  -> F_q
\newcommand{\bits}{\{0,1\}}
\newcommand{\bitstr}[1]{\{0,1\}^{#1}}         % \bitstr{n}

% ============================================================
% NOTATION — SCHEMES AND OPERATIONS
% ============================================================

\newcommand{\Gen}{\mathsf{Gen}}
\newcommand{\Enc}{\mathsf{Enc}}
\newcommand{\Dec}{\mathsf{Dec}}
\newcommand{\Sign}{\mathsf{Sign}}
\newcommand{\Vrfy}{\mathsf{Vrfy}}
\newcommand{\Mac}{\mathsf{Mac}}
\newcommand{\PRF}{\mathsf{F}}
\newcommand{\Hash}{\mathsf{H}}
\newcommand{\KDF}{\mathsf{KDF}}

\newcommand{\xor}{\oplus}
\newcommand{\concat}{\mathbin\Vert}
\newcommand{\getsr}{\stackrel{\scriptscriptstyle\$}{\leftarrow}} % x \getsr S

\newcommand{\secpar}{\lambda}
\newcommand{\adv}{\mathcal{A}}
\newcommand{\sect}[1]{\mathcal{#1}}
\newcommand{\msgspace}{\mathcal{M}}
\newcommand{\keyspace}{\mathcal{K}}
\newcommand{\ctspace}{\mathcal{C}}

\newcommand{\abs}[1]{\left\lvert #1 \right\rvert}
\newcommand{\ceil}[1]{\left\lceil #1 \right\rceil}
\newcommand{\floor}[1]{\left\lfloor #1 \right\rfloor}
\newcommand{\Prob}[1]{\Pr\!\left[ #1 \right]}

% ============================================================
% TIKZ — TWO-PARTY PROTOCOL DIAGRAMS
% ============================================================

\tikzset{
    party/.style        = {font=\bfseries},
    lifeline/.style     = {thick, rulegray, densely dashed},
    msg/.style          = {-{Stealth[length=2.6mm]}, thick, schoolred},
    msglabel/.style     = {font=\small, inner sep=2pt, fill=white},
    compute/.style      = {draw=rulegray, fill=black!3, rounded corners=2pt,
                           font=\small, inner sep=5pt, align=center}
}

% \begin{protocol}{caption}{label} ... tikz code ... \end{protocol}
\newenvironment{protocol}[2]
  {\def\protocolcaption{#1}\def\protocollabel{#2}%
   \begin{figure}[H]\centering\begin{tikzpicture}}
  {\end{tikzpicture}\caption{\protocolcaption}\label{\protocollabel}\end{figure}}
